\documentclass[pdflatex,sn-nature]{sn-jnl}
\usepackage[T1]{fontenc}
\usepackage{graphicx}%
\usepackage{multirow}%
\usepackage{amsmath,amssymb,amsfonts}%
\usepackage{amsthm}%
\usepackage{mathrsfs}%
\usepackage[title]{appendix}%
\usepackage{xcolor}%
\usepackage{textcomp}%
\usepackage{manyfoot}%
\usepackage{booktabs}%
\usepackage{algorithm}%
\usepackage{algorithmicx}%
\usepackage{algpseudocode}%
\usepackage{listings}%
%%%%

\begin{document}
		\title{Information Sheet}
	
	\maketitle
	
	\paragraph{What is the main claim of the paper? Why is this an important contribution to the machine learning/data mining research area?}
	
	We introduce a novel metric that quantifies the attribution of polarization to specific annotator groups. Current literature relies on annotation agreement metrics for such analysis, even though agreement metrics have been shown to not be the best formulation for such analysis. Our metric bypasses many common issues such as restrictions on single-annotation setups, and problems when the annotator subgroups have greatly mismatched sizes. Unlike most other such metrics, we also offer a test for statistical significance.
	
	\paragraph{What is the evidence provided to support claims? Be precise.}
	Our mathematical and statistical formulation, as well as the replication of many annotation patterns found in relevant literature. More specifically, we find substantiative differences between annotators of different ethnicity, as reported in~\cite{goyal2022your}, gender~\cite{goyal2022your, binns2017like, giorgi2025human, al2020identifying}, and religious beliefs~\cite{sachdeva2022assessing}.
	
	\paragraph{Report 3-5 most closely related contributions in the past 7 years (authored by researchers outside the authors’ research group) and briefly state the relation of the submission to them.}
	Annotator behavior has been linked to different sociodemographic, ideological attributes, as well as personal experiences in previous work~\cite{ kumar-et-al-2021, sachdeva2022assessing, goyal2022your, binns2017like}. However, all these studies analyzed the annotation using agreement metrics, which may be unreliable. We instead perform similar analysis, using a tool that is explicitly designed to be reliable for subjective tasks with mismatched group sizes, both of which are common assumptions for minority groups used in content moderation, toxicity, and hate speech detection.
	
	\paragraph{Specify 5 general keywords and 5 specific keywords describing the main research activity presented in the manuscript.}
	
	\begin{itemize}
		\item General: Polarization, annotation, socio-demographic, statistics, metric
		\item Specific: Statistical significance, quantification, toxicity, hate speech, content moderation
	\end{itemize}

	
	\paragraph{Declare any conflict of interest by reporting the email domains of all institutions with which the authors have an institutional conflict of interest. Authors have an institutional conflict of interest if they are currently employed or have been employed at this institution in the past three years, or if the authors have extensively collaborated with this institution within the past three years. Authors are also required to identify all Guest Editorial Board Members with whom the authors have a conflict of interest. Examples of conflicts of interest include co-authorship in the last five years, a colleague in the same institution within the last five years, advisor/student relationships, parentage relationships.}
	aueb.gr, athenarc.gr, ...
	
	\paragraph{Who are the most appropriate reviewers for the paper? Authors are required to suggest up to four candidate reviewers (especially if external to the Guest Editorial Board), including a brief motivation for each suggestion. Optionally, list up to four researchers/potential reviewers with competing interests that should not be considered for reviewers. PhD students cannot be suggested as potential reviewers. For each cut-off date, there is paper bidding.}
	...
	
	
	\bibliography{refs}
	
\end{document}